\cleardoublepage
\section{Solving One-Step Equations}

In this section, we'll learn how to solve equations that can be finished in just one step. We'll use \textbf{inverse operations} to isolate the variable. That means undoing what's being done to the variable. Always \textbf{do the same thing to both sides} of the equation to keep it balanced.

Solving equations is a lot like solving a riddle or puzzle. You're looking for the one number that makes the equation true. Use inverse operations to isolate the variable. That means undoing what's being done to the variable.

\subsection*{Example: Using Inverse Operations}

Let's look at some examples of solving one-step equations:

\textbf{Addition Example:}
\[
x + 4 = 9 \quad \text{Subtract 4 from both sides} \quad x = 5
\]

\textbf{Multiplication Example:}
\[
5x = 20 \quad \text{Divide both sides by 5} \quad x = 4
\]

Let's say you have \$20, and you know that it's 5 times the cost of one notebook. You can write $5x = 20$ and solve for $x$ to find out the price of one notebook. Algebra helps you work backward to find unknowns.

When solving equations, ask: "What's happening to the variable?" Then do the opposite. Keep both sides of the equation equal—whatever you do to one side, do to the other. Always plug your answer back in to make sure it works!

\cleardoublepage
\subsection{Short Question (English)}
Solve: $x - 9 = 15$

\fullpageimage{images/q2_3_1-eng.jpg}

\newpage
\subsection{Short Question (Korean)}
Solve: $x - 9 = 15$

\fullpageimage{images/q2_3_1-kor.jpg}

\newpage
\subsection{Short Question (Answer)}
Solve: $x - 9 = 15$

\textbf{Answer:} $x = 24$

\cleardoublepage
\subsection{Long Question (English)}
Solve: $6x = 54$

\fullpageimage{images/q2_3_2-eng.jpg}

\newpage
\subsection{Long Question (Korean)}
Solve: $6x = 54$

\fullpageimage{images/q2_3_2-kor.jpg}

\newpage
\subsection{Long Question (Answer)}
Solve: $6x = 54$

\textbf{Answer:} $x = 9$

\cleardoublepage
\subsection{Challenge Question (English)}
Solve: $\frac{x}{8} = 6$

\fullpageimage{images/q2_3_3-eng.jpg}

\newpage
\subsection{Challenge Question (Korean)}
Solve: $\frac{x}{8} = 6$

\fullpageimage{images/q2_3_3-kor.jpg}

\newpage
\subsection{Challenge Question (Answer)}
Solve: $\frac{x}{8} = 6$

\textbf{Answer:} Multiply both sides by 8: $x = 48$

